\documentclass[aspectratio=169]{beamer}
\input{header}
\coursecode{JKSSB}

\title{CPU vs GPU \& Microprocessor}
\subtitle{Lesson 06 --- Serial vs Parallel, Microprocessor vs Microcontroller}
\author{JKSSB Computer Awareness}
\date{\today}

\begin{document}
\frame{\titlepage}

\begin{frame}{Video Lesson 06}
  \centering
  \Large\href{https://youtu.be/Vaod5Ox7lzM}{Watch: CPU vs GPU: What's the Difference?}
\end{frame}

\begin{frame}{Serial vs Parallel Processing}
  \begin{itemize}
    \item \key{Serial processing}: one task is carried out at a time, in
      sequence.
    \item \key{Parallel processing}: several tasks or operations proceed at
      the same time.
    \item A few fast cores suit serial work; many cores working together suit
      parallel work.
    \item The office desktop (EX-01) runs several programs together, so it
      depends on parallel support.
  \end{itemize}
  \begin{block}{Exam cue}
    Serial = order (A, then B, then C). Parallel = together (A and B and C).
  \end{block}
\end{frame}

\begin{frame}{CPU: Few Powerful Cores}
  \begin{itemize}
    \item The CPU is the \key{Central Processing Unit}, the general-purpose
      processor.
    \item It has \key{few, powerful cores}, each able to run a sequential
      instruction stream quickly.
    \item Design goal: \key{low latency} --- finish one task with minimum
      delay.
    \item In EX-01, the CPU runs the operating system, MS Word, and
      spreadsheet logic step by step.
  \end{itemize}
  \begin{alertblock}{Remember}
    A CPU core is general-purpose: it can run any instruction, not only
    graphics.
  \end{alertblock}
\end{frame}

\begin{frame}{GPU: Many Simple Cores}
  \begin{itemize}
    \item GPU is the \key{Graphics Processing Unit}.
    \item It has \key{many simpler cores}, each performing a small amount of
      work.
    \item Design goal: \key{high throughput} --- a huge number of parallel
      operations per second.
    \item Natural fit: pixels and graphics, plus scientific and
      machine-learning workloads.
  \end{itemize}
  \begin{alertblock}{Trap}
    \bad{GPU cores are not general-purpose CPU cores.} They are many,
    simpler, and optimised for parallel work.
  \end{alertblock}
\end{frame}

\begin{frame}{CPU vs GPU at a Glance}
  \begin{center}
    \begin{tabular}{ll}
      \toprule
      \textbf{Feature} & \textbf{CPU vs GPU} \\
      \midrule
      Cores & CPU: few, powerful \quad GPU: many, simple \\
      Optimised for & CPU: latency \quad GPU: throughput \\
      Best at & CPU: sequential \quad GPU: parallel \\
      Example & CPU: desktop logic \quad GPU: pixels \\
      \bottomrule
    \end{tabular}
  \end{center}
  \vspace{0.3cm}
  \muted{Think: a few experts (CPU) vs a large team of simple workers (GPU).}
\end{frame}

\begin{frame}{Cores vs Threads and Hyper-Threading}
  \begin{columns}[T]
    \column{0.46\textwidth}
      \textbf{Core}
      \begin{itemize}
        \item A physical processing unit.
        \item A quad-core CPU has four such units.
      \end{itemize}
    \column{0.46\textwidth}
      \textbf{Thread}
      \begin{itemize}
        \item A stream of instructions a core executes.
        \item \key{Hyper-threading}: one core presents two logical threads.
      \end{itemize}
  \end{columns}
  \vspace{0.2cm}
  \begin{block}{Do not confuse}
    Two logical threads on one core are \bad{not} two physical cores.
  \end{block}
\end{frame}

\begin{frame}{Cache Depth: L1, L2, L3}
  \begin{itemize}
    \item Cache is small, fast memory held close to the core.
    \item \key{L1}: smallest and fastest, nearest the core.
    \item \key{L2}: larger and a little slower than L1.
    \item \key{L3}: largest and slowest of the three, shared across cores.
    \item Speed order: L1 faster than L2 faster than L3, then main memory.
  \end{itemize}
  \begin{block}{Exam cue}
    Smallest $\rightarrow$ fastest $\rightarrow$ closest:
    \key{L1\,$<$\,L2\,$<$\,L3}.
  \end{block}
\end{frame}

\begin{frame}{Clock Speed and Multi-Core}
  \begin{itemize}
    \item \key{Clock speed} is measured in \key{GHz} (gigahertz): billions of
      cycles per second.
    \item A higher clock speed generally means more work completed per second,
      all else equal.
    \item A \key{multi-core} processor has several cores on one chip, which
      enables parallel work.
    \item Clock speed alone does not decide real speed --- cores, cache,
      memory, and I/O all matter.
  \end{itemize}
\end{frame}

\begin{frame}{Microprocessor: A CPU on a Chip}
  \begin{itemize}
    \item A \key{microprocessor} is a \key{central processing unit implemented
      on a single chip}.
    \item It is a \key{general-purpose} processor: software decides what it
      does.
    \item It needs external memory and I/O devices to form a complete working
      system.
    \item The office desktop (EX-01) is built around a microprocessor.
  \end{itemize}
  \begin{block}{Definition}
    Microprocessor = the CPU, integrated on one chip.
  \end{block}
\end{frame}

\begin{frame}{Microcontroller: Integration on One Chip}
  \begin{itemize}
    \item A \key{microcontroller} adds more onto one chip: a microprocessor
      \key{plus} memory \key{plus} I/O peripherals.
    \item It forms a near-complete computer on one chip, used for
      \key{embedded/special-purpose} control.
    \item Everyday uses: washing machines, remote controls, and appliance or
      vehicle controllers.
    \item The key difference from a microprocessor is the \key{level of
      integration}, not clock speed.
  \end{itemize}
\end{frame}

\begin{frame}{Microprocessor vs Microcontroller}
  \begin{center}
    \begin{tabular}{ll}
      \toprule
      \textbf{Microprocessor} & \textbf{Microcontroller} \\
      \midrule
      CPU on a chip & CPU + memory + I/O on a chip \\
      General-purpose & Embedded / special-purpose \\
      Needs external parts & Self-contained system \\
      E.g. office desktop & E.g. appliance controller \\
      \bottomrule
    \end{tabular}
  \end{center}
  \begin{alertblock}{Trap}
    \bad{Do not call a microcontroller ``just a microprocessor''} --- it is a
    more integrated device.
  \end{alertblock}
\end{frame}

\begin{frame}{Diagnosing a ``Slow PC'' (EX-01)}
  \begin{itemize}
    \item EX-01: a quad-core 2.4\,GHz desktop with 16\,GB RAM and an SSD.
    \item Yet one program still runs slowly.
    \item More cores do not help a \key{single-threaded} program.
    \item The cause is usually a \key{bottleneck}: memory, disk I/O, or
      software that is not multithreaded.
  \end{itemize}
  \begin{alertblock}{Trap (TRAP-11)}
    \bad{A slow program is not automatically a ``slow CPU.''} Diagnose the
    stated bottleneck first.
  \end{alertblock}
\end{frame}

\begin{frame}{Lesson 06: Recall}
  \begin{itemize}
    \item CPU: few powerful, general-purpose cores optimised for latency.
    \item GPU: many simple cores optimised for parallel throughput.
    \item Serial = one at a time; parallel = many at once.
    \item Microprocessor = CPU on a chip; microcontroller = CPU + memory +
      I/O on one chip.
    \item Quad-core but slow $\rightarrow$ suspect a bottleneck, not always
      the CPU.
  \end{itemize}
\end{frame}
\end{document}
